%-*-latex-*-
%-*-latex-*-
\newcommand\COURSE{ciss350}
\newcommand\ASSESSMENT{q04}
\newcommand\ASSESSMENTTYPE{Quiz}
\newcommand\POINTS{\textwhite{xxx/xxx}}

\input{myquizpreamble}
\input{yliow}
\input{\COURSE}
%-*-latex-*-

%-*-latex-*-

%-*-latex-*-
\renewcommand\TITLE{\ASSESSMENTTYPE \ \ASSESSMENT}

\renewcommand\EMAIL{}
\renewcommand\AUTHOR{eagraham1@cougars.ccis.edu}

\textwidth=6in
\begin{document}
\topmatterthree

%------------------------------------------------------------------------------
\nextq
Write a function that accepts two C-strings \verb!x! and \verb!y! and
computes the longest common substring in \verb!x! and \verb!y!:
\begin{console}
void longest_common_substring(const char * x, const char * y, char * z);
\end{console}
The longest common substring is placed in \verb!z!.

For instance if \verb!x! is \verb!"hello world"! and \verb!y! is
\verb!"Hello world"!, then \verb!z! is \verb!"ello world"!.
When \verb!x! and \verb!y! are empty strings
or when \verb!x! is \verb!"abc"! and \verb!y! is \verb!"def"!, the longest
common substring is \verb!""!.
If there are two longest common substrings, the first one that appears
in \verb!x! is used.
For instance if \verb!x! is \verb@"hello!world."@ and
\verb!y! is \verb@"world.hello!"@, then there are two longest substrings
that appear in \verb!x! and \verb!y!: \verb@"hello!"@ and \verb!"world."!,
both of length 6. However \verb@"hello!"@ appears earlier in \verb!x!.
Therefore \verb!z! is \verb@"hello!"@.

Just use the simplest algorithm you can come up with.
You don't have to be super efficient.
The \lq\lq correct" algorithm is covered in CISS358.
You may use the \verb!strlen! function from the C-string library.
You can add helper functions.

Recall that if \verb!x! has type \verb!const char *!,
it means that \verb!x! points to a \verb!char! and cannot change
the value that \verb!x! points to.
If you have a function of the form
\verb!void f(const char *)!
then you can call \verb!f("abc")!.
However you cannot call \verb!g("abc")! if \verb!g! has a prototype
of \verb!void g(char *)!.
This is because something like \verb!"abc"! is considered a
constant C-string.

TURN PAGE

\newpage
\ANSWER
\begin{answercode}
void longest_common_substring(const char * x, const char * y, char * z)
{
  int x_len = strlen(x);
  int y_len = strlen(y);
  
  int best_len = 0;
  int best_start = 0;
  
  for(int i = 0; i < x_len; i++)
  {
      for(int j = 0; j < y_len; j++)
      {
          int k = 0;
          
          while(i + k < x_len &&
                i + k < y_len &&
                x[i + k] == y[j + k])
          {
              k++;
          }
          
          if (k > best_len)
          {
              best_len = k;
              best_start = i;
          }
      }
  }
  for(int i = 0; i < best_len; i++)
      z[i] = x[best_start + i];
  z[best_len] = '\0';
}
\end{answercode}

TURN PAGE FOR SPOILER ONLY IF YOU WANT SOME HELP

\newpage
\textsc{Spoiler}

Forget about the code.
Think first.
This problem involves comparing substrings of two strings
\verb!x! and \verb!y!.
This means that if you need to systematically iterate over all substrings
of \verb!x! and all the substrings of \verb!y!,
and then compare the two substrings:
\begin{Verbatim}[fontsize=\small,frame=single]
for s running over all substrings of x:
    for t running over all substrings of y:
        compare s and t
\end{Verbatim}
The question is then: What is a substring of \verb!x!?

A substring of \verb!x! can be described by
\verb!x!, \verb!start!, and \verb!end!
where \verb!start! and \verb!end! are two index values.
The substring is made up of the characters from index
\verb!start! up to but not including \verb!end!, i.e.,
\verb!x[start .. end - 1]!.

For instance if \verb!x! is \verb!"hello"!
and \verb!start! is 1 and \verb!end! is 4, then the
substring is \verb!"ell"!.
With the above idea you should be able to iterate over
all substrings of \verb!"hello"! and get
\verb!""!,
\verb!"h"!,
\verb!"he"!,
\verb!"hel"!,
\verb!"hell"!,
\verb!"hello"!,
\verb!""!,
\verb!"e"!,
\verb!"el"!,
\verb!"ell"!,
\verb!"ello"!,
\verb!""!,
\verb!"l"!,
\verb!"ll"!,
etc.
Believe it or not, you have seen this double for-loop before in my
notes on for-loops from CISS240 (the second set of notes on for-loops).
If you print the index values processed by this double for-loop
you might recognize it.

Once you can do the above, for each substring of \verb!x!,
you iterate over all substrings of
\verb!y!, and compare the two substrings.
There are some obvious optimizations of the above idea.
But even with these optimizations, the algorithm is still going to be extremely
slow.
The \lq\lq correct" algorithm for this very important algorithm
will be covered in CISS358.

Also, instead of using two index values for \verb!x!, you can use
two pointers.
In fact it's a good idea to use pointers just for the sake of reviewing
pointers.

Here's a test code:
\begin{Verbatim}[fontsize=\small,frame=single]
void test(const char * x, const char * y)
{
    char z[1024];
    std::cout << "Test [" << x << "] " << "[" << y << "]: ";
    longest_common_substring(x, y, z);
    std::cout << "[" << x << "] "
              << "[" << y << "] "
              << "[" << z << "]\n";
}

int main()
{
    test("abc", "def");
    test("", "def");
    test("abc", "");
    test("abc", "abc");
    test("abc", "Abc");
    test("Abc", "abc");
    test("abC", "abc");
    test("abc", "abC");
    test("abcd", "AbcD");
    test("AbcD", "abcd");
    test("abcd", "AAAAAAbcD");
    test("AAAAAAbcD", "abcd");
    test("abcd", "AAAAAAbcDDDDDDD");
    test("abcd", "AAAAAAbcDDDDDDD");
    test("AAAAAAbcDDDDDDD", "abcd");
    test("abcd", "AabAabAAabcDDDDDabcdDD");
    test("AabAabAAabcDDDDDabcdDD", "abcd");
    test("abcd", "AabcdAAAabcDDDDDabdDaD");
    test("abcd", "AAAAAAbcDDDDabDDD");
    test("abcd", "AAAAAAbcdDDDDabcDDD");
    test("AAAAAAbcdDDDDabcDDD", "!!!!!!!abcd@@@@@@@");

    return 0;
}
\end{Verbatim}

TURN PAGE FOR MORE SPOILER

\newpage
\textsc{More spoilers}

Here's a function that will be helpful for printing while debugging:
\begin{Verbatim}[fontsize=\small,frame=single]
// Print p[0..len-1] where len = q - p.
// "..............abcde...................."
//                ^    ^
//                |    |
//                p    q
void print(const char * p, const char * q)
{
    while (p < q)
    {
        std::cout << (*p);
        ++p;
    }
}
\end{Verbatim}
And here's a function that will compare two subarrays of characters:
\begin{Verbatim}[fontsize=\small,frame=single]
// Return true if r[0 .. len-1] is p[0 .. len-1] where len = s - r.
//
// "...abcde...................."
//     ^
//     |
//     p
// "..............abcde...................."
//                ^    ^
//                |    |
//                r    s
bool isprefix(const char * p, const char * r, const char * s)
{
    while (r < s)
    {
        if (*p != *r)
        {
            return false;
        }
        ++p; ++r;
    }
    return true;
}
\end{Verbatim}

%------------------------------------------------------------------------------
%-*-latex-*-
\newpage
\input{instructions}
\end{document}

